\documentclass{ufctex}

\ufcsetup{
  tipo = tese,
  impressao = frente-verso,
  capa = nao,
  brasao = nao,
  ies = {Universidade Federal do Ceará},
  programa-doutorado = {Programa de Pós-Graduação em Ciência da Computação},
  titulo-doutor = {Ciência da Computação},
  autor = {Autor Duplex Teste},
  titulo = {Validação Duplex dos Pré-textuais},
  local = {Fortaleza},
  ano = {2026},
  data-aprovacao = {19 de agosto de 2026},
  orientador = {Prof. Dr. Orientador Duplex},
  orientador-ies = {Universidade Federal do Ceará}
}

\begin{document}
\pretextual

\imprimirfolhaderosto
\imprimirerrata{tests/fixtures/pretextuais/errata.tex}
\imprimirfolhadeaprovacao
\imprimirdedicatoria{tests/fixtures/pretextuais/dedicatoria.tex}
\imprimiragradecimentos{tests/fixtures/pretextuais/agradecimentos.tex}
\imprimirepigrafe[longa]{tests/fixtures/pretextuais/epigrafe.tex}
\imprimirresumo{tests/fixtures/pretextuais/resumo.tex}
\imprimirabstract{tests/fixtures/pretextuais/abstract.tex}
\imprimirlistadeabreviaturasesiglas{tests/fixtures/pretextuais/abreviaturas.tex}
\imprimirlistadesimbolos{tests/fixtures/pretextuais/simbolos.tex}
\imprimirsumario

\textual
\section{Introdução duplex}
MARCADOR-TEXTUAL-DUPLEX.

\end{document}
